\chapter{Discussion}\label{ch:discussion}
% Looking back: consequences of choices made along the way, process and results.

\section{Assumptions About the Deployment Platform}\label{sec:deployment-assumptions}

Authentication was the first part of the system where a component behaved
differently in deployment than it did locally, and the difference was silent.

The Keycloak realm is kept in version control as \texttt{realm-prepmeup.json} so
the identity configuration can be recreated rather than clicked together by hand.
The container imports it at startup with \texttt{-{}-import-realm}, and locally
this worked: the file was made available to the container through a bind mount
declared in the compose file. Deployed on Coolify, the same declaration produced
a realm that did not exist. Requests to
\texttt{/realms/prepmeup/.well-known/openid-configuration} returned 404, while
the container reported a healthy startup and logged no error.

The cause was the relative path in the bind mount. Coolify resolves such paths
against its own working directory rather than against the cloned repository, so
Docker found no file at the source path. Rather than failing, it created an empty
directory at the mount target, which was confirmed on the server:

\begin{verbatim}
$ docker exec <keycloak-container> ls -la /opt/keycloak/data/import/
drwxr-xr-x 2 root root 4096 Sep 16 12:25 realm-prepmeup.json
\end{verbatim}

The directory bit in \texttt{drwxr-xr-x} shows that the realm file had become a
directory. Keycloak scanned an empty import folder, found nothing to import, and
started normally. Every individual component behaved as documented; the failure
lived in the assumption that a relative path means the same thing on both sides
of a deployment.

The fix was to stop depending on the platform's path handling and copy the realm
into a thin image layered on the Keycloak base image, so the file is part of the
build artefact rather than something mounted at run time. A second constraint
surfaced during the same fix: \texttt{-{}-import-realm} only imports against an
uninitialised database, so the existing Keycloak volume had to be discarded
before the corrected configuration took effect. That cost nothing at this stage,
but it would have been a destructive operation with real users in the system, and
it is the reason the realm file is treated as the source of truth rather than the
running database.

Two lessons carry forward. The first is that a deployment platform is part of the
system under test: a configuration that works under \texttt{docker compose} on a
developer machine has not been shown to work in production, because the platform
resolves paths, networks and volumes on its own terms. The second is that silent
success is more expensive to diagnose than a crash. Keycloak treated an empty
import directory as an empty task rather than a misconfiguration, so the defect
only became visible one layer up, when the API could not be pointed at an issuer
that did not exist.

This informed how JWT validation was then configured. The API pins the expected
issuer explicitly instead of inferring it from the discovery document, because
the two addresses legitimately differ in deployment: discovery is fetched over
the internal container network, while the tokens clients present carry the public
issuer. Making that distinction explicit in configuration means the mismatch
becomes a validation error with an obvious cause, rather than another failure
that depends on where the code happens to be running.
